\textbf{Chi tiết cài đặt}

Thí sinh cần cài đặt hàm sau:

\begin{lstlisting}
void solve(int N);
\end{lstlisting}

\begin{itemize}
   \item $N$: Số đỉnh của cây bạn được cho trước.
   \item Hàm này cần mô tả cách chương trình của bạn hoạt động để thực hiện các yêu cầu của bài toán đã được mô tả ở trên.
   \item Hàm này được gọi $T$ lần trong một test, mỗi lần gọi tương ứng với một trường hợp thử nghiệm.
\end{itemize}

Trong hàm \texttt{solve()}, thí sinh được phép gọi các hàm sau:

\begin{lstlisting}
void print(int i);
\end{lstlisting}

\begin{itemize}
   \item $i$: Chỉ số của ô nhớ cần in ra. Cần phải đảm bảo $0\le i\le 1999$.
   \item Hàm này sẽ mô phỏng việc in ra giá trị của ô nhớ thứ $i$ ra màn hình.
\end{itemize}

\begin{lstlisting}
void sub(int i, int j, int k);
\end{lstlisting}

\begin{itemize}
   \item $i,j,k$: Chỉ số của các ô nhớ cần thiết để thực hiện thao tác cập nhật. Cần phải đảm bảo $0\le i,j,k\le 1999$ và các giá trị này không nhất thiết phải đôi một phân biệt.
   \item Hàm này sẽ thực hiện cập nhật $A_k:=A_i-A_j$.
\end{itemize}

\begin{lstlisting}
void mul(int i, int j, int k);
\end{lstlisting}

\begin{itemize}
   \item $i,j,k$: Chỉ số của các ô nhớ cần thiết để thực hiện thao tác cập nhật. Cần phải đảm bảo $0\le i,j,k\le 1999$ và các giá trị này không nhất thiết phải đôi một phân biệt.
   \item Hàm này sẽ thực hiện cập nhật $A_k:=A_i\times A_j$.
\end{itemize}

\begin{lstlisting}
void cmp(int i, int j, int k);
\end{lstlisting}

\begin{itemize}
   \item $i,j,k$: Chỉ số của các ô nhớ cần thiết để thực hiện thao tác cập nhật. Cần phải đảm bảo $0\le i,j,k\le 1999$ và các giá trị này không nhất thiết phải đôi một phân biệt.
   \item Hàm này sẽ thực hiện cập nhật $A_k:=1$ nếu $A_i=A_j$ và $A_k:=0$ trong trường hợp ngược lại.
\end{itemize}

\textbf{Lưu ý:}
\begin{itemize}
   \item Chương trình của bạn cần phải khai báo thư viện \texttt{"vmachinelib.h"} ở đầu.
   \item Với mỗi trường hợp thử nghiệm tương ứng với một lần gọi hàm \texttt{solve()}, chương trình của bạn được phép thực hiện tối đa $4\times 10^6$ lệnh.
   \item Trong chương trình, thí sinh được phép cài đặt các biến, các hàm toàn cục, nhưng không được phép có hàm \texttt{main()}.
   \item Chương trình của bạn không được tương tác trực tiếp với đầu vào chuẩn (stdin) và đầu ra chuẩn (stdout).
\end{itemize}

\textbf{Ràng buộc}

\begin{itemize}
   \item $T \le 50$
   \item $N \le 200$
   \item Độ dài của dãy $A$ (kích thước của khe nhớ) luôn bằng $2000$
\end{itemize}